\section{Misura utilizzata}
Poiché Linux e Windows offrono statistiche diverse riguardanti la memoria allocata per un processo, è stato
necessario utilizzare due metriche di memoria diverse:
\begin{itemize}
    \item \textbf{Linux}: Per Linux abbiamo utilizzato la misura \textbf{Resident Set Size} (RSS), la quale indica la memoria effettivamente occupata 
    \item \textbf{Windows}:  
\end{itemize}





\section{Profiling della memoria}
Per effettuare le misurazioni sullo spazio occupato in memoria dei due programmi sviluppati,
è stato necessario scrivere due ulteriori programmi distinti, uno per Windows e uno per Linux, per il campionamento della memoria occupata da un
determinato processo.
Questa necessità nasce da problematiche legate sia a Matlab che a C++.
Per quanto riguarda Matlab, la pressoché inesistente documentazione sul profiler integrato nell'ambiente di sviluppo \cite{mathworks_profiling}
rende particolarmente arduo comprendere appieno la misurazione ottenuta. 
Per C++, invece, non abbiamo trovato un profiler agnostico al sistema operativo che indicasse la quantità di memoria utilizzata per ogni chiamata
ad una certa funzione (nel nostro caso, la funzione \texttt{calculate\_benchmark()}).
Di seguito, riportiamo i due programmi scritti. I linguaggi utilizzati sono \textbf{Powershell} per Windows e \textbf{Bash} per Linux.

\lstinputlisting[
    language=Powershell,
    style = PowerShell-editor,
    caption = Sampler scritto per Windows in Powershell
]{../profilers/mem_sampler.ps1}
\clearpage
\lstinputlisting[
    language=Bash,
    style = Bash-editor,
    caption = Sampler scritto per Linux in Powershell
]{../profilers/mem_sampler.sh}